\section{An interpolation algorithm: technique overview}\label{sec:technical-overview}
Sections~\ref{sec:technical-overview}--\ref{sec:interpolation-return}
give an algorithm using polynomial interpolation over $\mathbb Q$ with exponent $O_k(\rho(n))$.
It uses the same parity tests as the main theorem and a different
method for counting their common zeros. This method separates the
summation routine from the construction of the replacement forms:
Section~\ref{sec:counting} gives the routine for arbitrary polynomial
factors over a commutative $\mathbb Q$-algebra, with the replacement
forms and measurements supplied. Its development is independent of
the contraction recursion in Section~\ref{app:tensor}.
We first illustrate the routine by counting satisfying assignments
of a $3$-CNF in which each variable occurs in at most three clauses. The product of its clause
polynomials has degree at most $3n$, so we can use interpolation to
recover the exact count from sums over small groups of variables.
Section~\ref{subsec:preview-general-sat} then uses disjoint
sparsification and variable copying to count the satisfying assignments
of an arbitrary $k$-CNF in $(n+m)^{O_k(1)}2^{o(n)}$ time for every
fixed $k\ge2$. Taking $k=3$ already refutes ETH.
Section~\ref{subsec:preview-full-case} explains how counting on the
original variables gives the exponent $O_k(\rho(n))$ of
Proposition~\ref{prop:interpolation-bound}.

\subsection{The bounded-occurrence setting}\label{subsec:preview-setting}
Let $F$ have $n$ declared Boolean variables. Each clause contains
at most three literals, and each variable appears in at most
three clauses, counting positive and negative occurrences together.
Our task is to count the satisfying assignments of $F$ exactly.

We remove tautological clauses and repeated literals within a
clause. An empty clause gives count zero immediately.
These operations preserve the occurrence bound. For each remaining
clause $e$, let $I_e\subseteq[n]$ be its variable set.
For each variable that occurs in no clause, add a factor equal to
one with scope $\{v\}$. Such a factor leaves the count unchanged
and ensures that every variable has an assigned position in the
construction.
Let $\mathcal F$ index the clause polynomials and these constant-one factors. Then
\begin{equation}\label{eq:overview-incidences}
 |I_e|\le3,\qquad
 \#\{e\in\mathcal F:v\in I_e\}\le3,\qquad
 n\le M:=\sum_{e\in\mathcal F}|I_e|\le3n.
\end{equation}
All sums below still range over the original $n$ Boolean variables.

\subsection{A direct polynomial for the count}\label{subsec:preview-encoding}
Write $L_{e,v}(X_v)=X_v$ for a positive literal and
$L_{e,v}(X_v)=1-X_v$ for a negative literal. The satisfaction
polynomial of clause $e$ is
\[
 \Phi_e((X_v)_{v\in I_e})
 =1-\prod_{v\in I_e}(1-L_{e,v}(X_v)).
\]
For example, $x_a\vee\neg x_b\vee x_c$ gives
\[
 \Phi_e(X_a,X_b,X_c)=1-(1-X_a)X_b(1-X_c).
\]
On Boolean inputs this polynomial is one when the clause is
satisfied and zero otherwise. For a constant-one factor $e$, put
$\Phi_e=1$. Therefore
\begin{equation}\label{eq:overview-direct-count}
 \#F=\sum_{x\in\{0,1\}^n}\prod_{e\in\mathcal F}
                   \Phi_e((x_v)_{v\in I_e}).
\end{equation}
The product has total degree at most $M\le3n$.
This linear degree bound is the useful consequence of the
three-occurrence assumption.

Directly evaluating \eqref{eq:overview-direct-count} would be
exhaustive search. If the variables admitted a partition into groups
$S_1,\ldots,S_r$, each of size at most $2n/r$, with every clause
contained in one group, we could instead count within groups.
Writing $N_g$ for the number of assignments to the variables in $S_g$
that satisfy that group's clauses, independence gives
\[
 \#F=\prod_{g=1}^r N_g.
\]
The local enumerations examine at most
$\sum_{g=1}^r2^{|S_g|}\le r\,2^{2n/r}$ assignments,
with polynomial work per assignment.
An arbitrary formula need not admit such a partition. We will replace
the clause inputs so that groups interact only through a small
collection of shared linear summaries. Fixing those summaries will
again let us sum over each group separately.

\subsection{A polynomial family with easier evaluation points}
\label{subsec:preview-deformation}
We introduce a parameter $T$. At $T=0$, the replacements must
restore the original factor inputs. At selected nonzero points
$T=1,\ldots,Q$, they should let us separate the groups by fixing
shared linear summaries. Here $Q$ is the number of selected points,
and $R$ will bound the degree of the replacement coefficients.
We choose both parameters later.

Write $\vec{X}=(X_1,\ldots,X_n)$ for the vector of formal variables.
For each occurrence $v\in I_e$, use the polynomial coefficient vector
$\vec{P}_{e,v}(T)=(p_{e,v,1}(T),\ldots,p_{e,v,n}(T))$ to replace
$X_v$ in factor $e$ by the scalar inner product
\[
 \langle \vec{P}_{e,v}(T),\vec{X}\rangle
 =\sum_{i=1}^n p_{e,v,i}(T)X_i.
\]
We call this expression the \emph{replacement linear form} for
the occurrence of variable $v$ in factor $e$. It replaces the
input $X_v$ in the factor $\Phi_e$.
We require
\[
 \langle \vec{P}_{e,v}(0),\vec{X}\rangle=X_v,
 \qquad \deg \vec{P}_{e,v}\le R.
\]
Different factors may use different replacement linear forms
for the same variable.
Define
\begin{equation}\label{eq:simple-counting-polynomial}
 H(T)=\sum_{x\in\{0,1\}^n}\prod_{e\in\mathcal F}
       \Phi_e\bigl((\langle \vec{P}_{e,v}(T), x\rangle)_{v\in I_e}\bigr)
       \in\mathbb Q[T].
\end{equation}
At zero every replacement restores the original variable, while
the degree of each factor grows by at most a factor of $R$. Hence
\begin{equation}\label{eq:overview-value-degree}
 H(0)=\#F,\qquad \deg H\le MR\le3nR.
\end{equation}
We will recover $H(0)$ from Taylor coefficients at $t=1,\ldots,Q$.
The degree bound means that $O(1+nR/Q)$ coefficients per point
suffice, as Section~\ref{subsec:preview-interpolation} explains.
Our goal is therefore to make $R$ small relative to $Q$: many
points with low-degree replacements require fewer coefficients
at each point.

At a nonzero $t$, the value $\langle\vec{P}_{e,v}(t),x\rangle$
of a replacement linear form need not lie in $\{0,1\}$, even
when $x$ is Boolean. The transformed factors
therefore need not act as satisfaction indicators.

\subsection{Separating the sum by fixing measurements}
\label{subsec:preview-separation}
Fix $t\in[Q]$. We are given the coefficient vectors
$\vec{P}_{e,v}(t)$, whose entries are now constants. We explain
what structure lets us compute $H(t)$ quickly, and how to use it.
Throughout this subsection, write $S_g=S_{t,g}$ and $C_g=C_{t,g}$.

\begin{definition}[Partitions and measurements]
\label{def:overview-measurement-structure}
For the given coefficient vectors $\vec{P}_{e,v}(t)$ and parameters $(r,\tau,D)$,
the groups $S_g$, factor assignments $\sigma_e(t)$,
and measurement matrices $C_g$ satisfy the following four
conditions:
\begin{enumerate}
\item \emph{Small groups.} Partition the variables into disjoint
groups $[n]=S_1\sqcup\cdots\sqcup S_r$, with $|S_g|\le2n/r$.
\item \emph{One assigned group per factor.} Assign each factor
$e$ to one group $\sigma_e(t)\in[r]$. Condition~4 must hold with
this same group for every replacement linear form of factor $e$.
\item \emph{Few bounded measurements.} Each group $g$ has an
integer matrix $C_g$ with columns indexed by $S_g$ and $c_g$ rows.
Each row specifies one linear measurement of the variables in
$S_g$. There are at most $2\tau n$ rows in total, and the sum of
absolute values of the coefficients in each row is at most $D$.
For an assignment $x$, write
\[
 x_g=x|_{S_g},\qquad y_g=C_gx_g,\qquad
 \vec{y}=(y_1,\ldots,y_r).
\]
\item \emph{Dependence on other groups only through measurements.}
Suppose factor $e$ is assigned to group $g=\sigma_e(t)$.
Each replacement linear form of this factor may depend freely on the
variables in $S_g$. Its dependence on any other group $S_h$ must
be a linear combination of that group's measurements $y_h$.
Precisely, for every $v\in I_e$, there must be coefficient vectors
$\vec{\lambda}_{e,v,h}\in\mathbb Q^{c_h}$, for $h\ne g$, such
that, for every $x\in\mathbb Q^n$,
\begin{equation}\label{eq:outline-separation}
 \langle\vec{P}_{e,v}(t),x\rangle
 =\langle\vec{P}_{e,v}(t)|_{S_g},x_g\rangle
   +\sum_{h\ne g}\langle\vec{\lambda}_{e,v,h},y_h\rangle.
\end{equation}
Equivalently, $\vec{P}_{e,v}(t)|_{S_h}$ lies in the row space
of $C_h$ for every $h\ne g$.
\end{enumerate}
\end{definition}
The matrices $C_h$ are shared by all factors. A factor's assigned
group need not contain its original variables. Once $\vec{y}$
is fixed, however, the entire transformed factor
depends only on the variables in its assigned group.

Find the coefficient vectors $\vec\lambda_{e,v,h}$ by solving the
linear systems guaranteed by condition~4. For a measurement tuple
$\vec y=(y_1,\ldots,y_r)$, put
\[
 \mathcal X_{t,g}(y_g)
 =\{u\in\{0,1\}^{S_g}:C_gu=y_g\},\qquad
 \mathcal Y=\prod_{g=1}^r\{-D,\ldots,D\}^{c_g}.
\]
For $u\in\mathcal X_{t,g}(y_g)$, use \eqref{eq:outline-separation}
to evaluate the transformed factors assigned to $g$ from $u$ and
$\vec y$, and let $A_{t,g}(u;\vec y)$ be their product. It is one if
no factor is assigned to $g$. Each measurement of a Boolean
assignment is an integer in $[-D,D]$. Hence each full Boolean
assignment has one measurement tuple in $\mathcal Y$ and one
restriction to each group, and its product of factors is
the product of the values $A_{t,g}$ at these restrictions.
Conversely, because $S_1,\ldots,S_r$ partition $[n]$, any choice of
$u_g\in\mathcal X_{t,g}(y_g)$ for every $g$ combines into exactly one
assignment $x$ with $x_g=u_g$ for all $g$, and the measurement tuple
of this $x$ is $\vec y$. Thus, for each $\vec y\in\mathcal Y$,
restriction to the groups is a bijection from the assignments with
measurement tuple $\vec y$ onto $\prod_{g=1}^r\mathcal X_{t,g}(y_g)$.
Distributing the product over these restrictions therefore gives
\begin{equation}\label{eq:outline-separated-sum}
 H(t)=\sum_{\vec{y}\in\mathcal Y}\prod_{g=1}^r
       \left(\sum_{u\in\mathcal X_{t,g}(y_g)}
                         A_{t,g}(u;\vec{y})\right).
\end{equation}
An empty measurement fiber contributes a zero group sum.

Evaluate this identity one measurement tuple at a time. For each
group, scan its Boolean assignments, retaining those with $C_gu=y_g$,
and accumulate the inner sum. Multiply the group sums and add the
product to the running total. There are at most
$(2D+1)^{2\tau n}$ measurement tuples and at most
$\sum_g2^{|S_g|}\le r2^{2n/r}$ assignments to scan per tuple.
The time is consequently
\[
 (2D+1)^{2\tau n}2^{2n/r}
\]
up to polynomial factors. Processing the groups sequentially retains
only the current measurement tuple, assignment, and running sums and
products; Lemma~\ref{lem:measurement-separation} gives the general
routine.

For the parameter choices in Section~\ref{subsec:preview-parameters},
we can construct replacement forms of degree at most $R$ that
restore the original variables at zero and admit partitions, factor assignments, and measurement matrices
with parameters $(r,\tau,D)$ at every $t\in[Q]$.
Section~\ref{subsec:preview-auxiliary} gives the construction.
We first explain how these properties let us recover the count
efficiently.

\subsection{Recovering the count from Taylor coefficients}
\label{subsec:preview-interpolation}
Our goal is to recover $H(0)=\#F$ from information at the points
$t=1,\ldots,Q$. Since $\deg H\le3nR$, the $Q$ values $H(t)$ alone
may not determine $H$. We instead compute several Taylor
coefficients at each point. Choose
\[
 q=\left\lfloor\frac{3nR}{Q}\right\rfloor,\qquad J=q+1.
\]
Suppose we know the coefficients of $z^0,\ldots,z^q$ in $H(t+z)$
for every $t\in[Q]$. These data determine $H$ uniquely: if two
polynomials of degree at most $3nR$ have the same data, their
difference is divisible by $(T-t)^J$ for every $t\in[Q]$, hence by
$\prod_{t=1}^Q(T-t)^J$. This product has degree $QJ>3nR$, so the
difference must be zero. Each Taylor coefficient gives a linear
equation in the coefficients of $H$. Solving this rational linear
system recovers $H$, and its constant coefficient is the desired
count $H(0)$.

It remains to compute these Taylor coefficients. Fix $t\in[Q]$;
we want $H(t+z)$ through degree $q$ in $z$. The row-space condition controls
the replacement forms at $t$, but need not control their changes
from $t$ to $t+z$. We first give those changes a separate input,
so that fixing this input preserves the separation of groups.

\smallskip\noindent
\textbf{Keeping the shifted inputs separate.}
Write
\[
 \Delta\vec P_{e,v}(z)=\vec P_{e,v}(t+z)-\vec P_{e,v}(t).
\]
Keep $x$ as the Boolean assignment being summed over, and introduce
a separate formal vector $\vec X$ for the changes. Define
\begin{equation}\label{eq:outline-taylor-product}
 \mathcal G_x(\vec X,z)=
 \prod_{e\in\mathcal F}\Phi_e\bigl(
       (\langle\vec P_{e,v}(t),x\rangle+
        \langle\Delta\vec P_{e,v}(z),\vec X\rangle)_{v\in I_e}\bigr).
\end{equation}
Here $x$ indexes the polynomial, while $\vec X$ and $z$ are its
formal variables. Setting $\vec X=x$ restores each shifted input
$\langle\vec P_{e,v}(t+z),x\rangle$, so
\[
 H(t+z)=\sum_{x\in\{0,1\}^n}\mathcal G_x(x,z).
\]
On the other hand, fixing $\vec X=\mathbf1_U$ for some
$U\subseteq[n]$ makes every correction
$\langle\Delta\vec P_{e,v}(z),\mathbf1_U\rangle$ independent of $x$.
All dependence on $x$ then remains through
$\langle\vec P_{e,v}(t),x\rangle$, where the row-space condition applies.
Our task is therefore to recover the sum with $\vec X=x$
from sums with $\vec X$ fixed.

The degree bound makes this possible with few fixed substitutions.
Every entry of $\Delta\vec P_{e,v}(z)$ is divisible by $z$, so
each occurrence of an $X_i$ in an expanded term contributes at
least one factor of $z$. Thus
\[
 \deg_{\vec X}([z^j]\mathcal G_x(\vec X,z))\le j.
\]
Thus, after discarding powers of $z$ above $q$, every monomial
involves at most $q$ distinct variables from $\vec X$.

\smallskip\noindent
\textbf{Recovering the shifted sum by inclusion--exclusion.}
Work modulo $z^J$, so that only the coefficients through $z^q$
are retained. Replace every positive power of $X_i$ in
$\mathcal G_x$ by $X_i$. This preserves its values on Boolean
inputs. The degree bound above gives, for Boolean $\vec X$,
\[
 \mathcal G_x(\vec X,z)\equiv
 \sum_{\substack{E\subseteq[n]\\|E|\le q}}
       a_E(x,z)\prod_{i\in E}X_i\pmod{z^J},
\]
where $a_E(x,z)$ is the coefficient of the indicated monomial
after this replacement. Evaluating at $\mathbf1_U$, the vector
with ones precisely in $U$, sums the coefficients indexed by
subsets of $U$. Inclusion--exclusion extracts each one:
\[
 a_E(x,z)\equiv
 \sum_{U\subseteq E}(-1)^{|E|-|U|}
                 \mathcal G_x(\mathbf1_U,z)\pmod{z^J}.
\]
When we set $\vec X=x$, the monomial $\prod_{i\in E}x_i$
equals one exactly when all the bits in $E$ are one.
Its contribution to the sum over $x$ is therefore
$\sum_{x:\,x_i=1\ (i\in E)}a_E(x,z)$, with $x$ Boolean.
Substituting the coefficient formula above reduces this to sums
with the formal input fixed.

\smallskip\noindent
\textbf{Computing the restricted sums.}
For $U\subseteq E\subseteq[n]$ with $|E|\le q$, define
\begin{equation}\label{eq:outline-restricted-taylor-sum}
 \mathcal Z_{t,E,U}(z)=
 \sum_{\substack{x\in\{0,1\}^n\\x_i=1\ (i\in E)}}
                    \mathcal G_x(\mathbf1_U,z)
       \quad\text{in }\mathbb Q[z]/(z^J).
\end{equation}
Here $U$ fixes the formal input $\vec X=\mathbf1_U$, while $E$
fixes bits of the summation variable $x$ to one.
Summing the monomial contributions gives
\begin{equation}\label{eq:outline-taylor-recovery}
 H(t+z)\equiv
 \sum_{\substack{E\subseteq[n]\\|E|\le q}}
       \sum_{U\subseteq E}(-1)^{|E|-|U|}
                         \mathcal Z_{t,E,U}(z)\pmod{z^J}.
\end{equation}
This recovers all $J=q+1$ Taylor coefficients together.
Lemma~\ref{lem:boolean-differences} states the general identity
used here.

Each sum $\mathcal Z_{t,E,U}$ can be evaluated by the algorithm
of Section~\ref{subsec:preview-separation}. With
$\vec X=\mathbf1_U$ fixed, each correction
$\langle\Delta\vec P_{e,v}(z),\mathbf1_U\rangle$ is independent of $x$.
All dependence on $x$ comes through the forms
$\langle\vec P_{e,v}(t),x\rangle$, so fixing $\vec y$ still separates
the factors. The restrictions $x_i=1$ are local to each
group. The separation identity uses only addition, multiplication,
and distributivity, so it also applies to these truncated polynomial
values. Each factor is evaluated after substituting
$\vec X=\mathbf1_U$, using only univariate polynomials in $z$
of degree less than $J$ and discarding higher powers after each
multiplication. The multivariate expansion and its coefficients
$a_E(x,z)$ are used only to justify the identity; the algorithm
never constructs them.

Repeating this calculation at $t=1,\ldots,Q$ provides all the
Taylor coefficients needed for the interpolation above, and
therefore recovers $H(0)=\#F$.

\subsection{The three costs and the running time}
\label{subsec:preview-parameters}
Jumping ahead, Section~\ref{subsec:preview-auxiliary} will construct
replacement forms that restore the original variables at zero,
have degree at most $R$, and admit partitions, factor assignments, and measurement matrices with parameters $(r,\tau,D)$ at every
$t\in[Q]$. For sufficiently large $n$, the construction takes
polynomial time and achieves the following parameter bounds:
\begin{equation}\label{eq:overview-parameter-bounds}
 \begin{aligned}
 r&=\Theta\!\left(\frac{\sqrt{\log n}}{(\log\log n)^2}\right),
 &Q&=\Theta\!\left(\frac{\log n}{\log\log n}\right),\\
 R&=\Theta(\sqrt{\log n}),
 &\tau&=\Theta\!\left(\frac1{r\log n}\right),\qquad
 D=r^{2Q}\le(2n)^{1/128}.
 \end{aligned}
\end{equation}
We first show how these bounds give a fast counting algorithm;
the next subsection will justify that they can all be achieved
simultaneously.

Recall that $q=J-1=\lfloor3nR/Q\rfloor$. For each evaluation
point $t$, Section~\ref{subsec:preview-interpolation} computes one
restricted sum for each pair $U\subseteq E$ with $|E|\le q$.
There are $\binom nj$ choices of an $E$ of size $j$, and then
$2^j$ choices of $U$. Each restricted sum enumerates measurement
tuples and assignments within individual groups. Thus the three
costs multiply to give, up to polynomial factors,
\begin{equation}\label{eq:overview-three-costs}
 \underbrace{\sum_{j=0}^{\min(q,n)}2^j\binom nj}_{\text{Taylor data}}
 \ \cdot\
 \underbrace{(2D+1)^{2\tau n}}_{\text{measurements}}
 \ \cdot\
 \underbrace{2^{2n/r}}_{\text{one group}}.
\end{equation}
To bound the first cost, note that $1\le q\le n/2$ for
sufficiently large $n$, since $R/Q\to0$. Using
$\binom nj\le(en/j)^j$ and the fact that $(2en/j)^j$ increases
for $1\le j\le n$, we obtain
\[
 \sum_{j=0}^q2^j\binom nj
 \le(q+1)\left(\frac{2en}{q}\right)^q.
\]
Since $q=\Theta(nR/Q)$, the logarithm of this cost is
$O(n(R/Q)\log(2Q/R))$. Taking logarithms of all three costs in
\eqref{eq:overview-three-costs} therefore gives
\[
 O\!\left(n\left[
       \frac RQ\log\frac{2Q}{R}
       +\tau\log(2D+1)+\frac1r\right]\right).
\]
Each term is small under \eqref{eq:overview-parameter-bounds}:
\[
 \begin{aligned}
 \frac RQ\log\frac{2Q}{R}
   &=O\!\left(\frac{\log\log n}{\sqrt{\log n}}
                      \cdot\log\log n\right),\\
 \tau\log(2D+1)&=O\!\left(\frac1{r\log n}\cdot\log n\right)
                  =O(1/r),\\
 \frac1r&=O\!\left(\frac{(\log\log n)^2}{\sqrt{\log n}}\right).
 \end{aligned}
\]
Thus the bounded-occurrence instance can be counted in
\begin{equation}\label{eq:overview-bounded-time}
 n^{O(1)}2^{O(n(\log\log n)^2/\sqrt{\log n})}
\end{equation}
bit operations and polynomial space. The finitely many smaller
inputs are handled by exhaustive search.

The number of points $Q$, the number of groups $r$, and the
truncation order $J=O(1+nR/Q)$ are polynomial in $n$. Hence the
loops over points and groups, arithmetic on polynomials of degree
less than $J$, and the final interpolation add only polynomially
many operations per iteration of the three enumerations above.

For the bit-complexity and space bounds, see
Appendix~\ref{subsec:overview-arithmetic}.

\subsection{Constructing the low-degree replacement linear forms}
\label{subsec:preview-auxiliary}
Consider one occurrence of $X_v$ in a factor. At $T=0$ its
replacement must equal $X_v$. At each $t=1,\ldots,Q$, the factor
must be able to evaluate that same replacement using only its
assigned group and a few shared measurements. Choosing a different
linear form at each point and interpolating would generally cost
degree $Q$. To keep Taylor recovery cheap, we need degree $R$
much smaller than $Q$.

Suppose a term uses a variable outside the factor's assigned group
only at $t=2$. Multiplying it by $1-T/2$ preserves its value at zero
and removes it at $t=2$. More generally, one such factor per
incompatible point suffices. We therefore seek a decomposition of
each input into a part determined by measurements and terms that
are incompatible at at most $R-1$ points. We describe the mechanism
here; Section~\ref{sec:construction} proves the construction for
general factor scopes.

\paragraph{The auxiliary decomposition.}
An occurrence has a variable-group choice and a factor-group choice
at every point. Index all possible choices by
\[
 \mathcal U=([r]\times[r])^Q,\qquad D=r^{2Q},\qquad
 u=((a_u(t),d_u(t)))_{t\in[Q]}.
\]
For placeholders $\vec\zeta=(\zeta_w)_{w\in\mathcal U}$, the
replacement for $\zeta_u$ may freely use $\zeta_w$ at point $t$
when $a_w(t)=d_u(t)$. Put
\[
 E_{u,w}=\{t\in[Q]:a_w(t)\ne d_u(t)\}.
\]
We need a rational matrix $M_0$ and an integer measurement matrix
$C\in\{0,1,-1\}^{c\times D}$ satisfying
\begin{equation}\label{eq:outline-matrix-conditions}
 \begin{gathered}
 c\le D\tau/(3r),\qquad
 M_0\vec\xi=\vec\xi\quad(\vec\xi\in\ker C),\\
 M_0[u,w]\ne0\quad\Longrightarrow\quad |E_{u,w}|\le R-1.
 \end{gathered}
\end{equation}
The kernel condition makes $(I_D-M_0)\vec\zeta$ a linear function
of $C\vec\zeta$. Thus the identity
\[
 \zeta_u=((I_D-M_0)\vec\zeta)_u+
                   \sum_w M_0[u,w]\zeta_w
\]
has the desired form. Define its replacement by
\begin{equation}\label{eq:outline-auxiliary-row}
 \langle\vec P_u(T),\vec\zeta\rangle
 =((I_D-M_0)\vec\zeta)_u+
   \sum_{w:M_0[u,w]\ne0}M_0[u,w]
       \prod_{t\in E_{u,w}}(1-T/t)\,\zeta_w.
\end{equation}
It equals $\zeta_u$ at zero and has degree at most $R-1$.
At point $t$, every surviving term outside the measured correction
belongs to group $d_u(t)$. This is the replacement property proved
in Lemma~\ref{lem:swaps-forms}.

\paragraph{Why a matrix can satisfy both conditions.}
Let $\Pi_s$ be the permutation matrix that
swaps the two entries of pair $s$ in every index, leaving all
other pairs unchanged.
The commuting swaps give
\[
 V_0=\prod_{s=1}^Q\Pi_s,\qquad
 W_0=\tfrac12\sum_{s=1}^Q(I_D-\Pi_s),\qquad M_0=V_0f(W_0).
\]
The complete swap $V_0$ connects matching groups at every point.
A polynomial $f$ of degree at most $R-1$ can undo at most $R-1$
swaps in any term of its expansion. Hence its nonzero entries have
at most $R-1$ incompatible points
(Lemma~\ref{lem:swaps-support}).

To make $M_0$ fix most directions, use the tensor products of the
symmetric and antisymmetric eigenvectors of a single swap. A basis
vector with $j$ antisymmetric factors has eigenvalues $j$ for $W_0$
and $(-1)^j$ for $V_0$. For a uniformly chosen basis vector,
$j$ has distribution $W\sim\operatorname{Bin}(Q,(r-1)/(2r))$.
Interpolate $f(j)=(-1)^j$ on an interval $\mathcal I$ of $R$
consecutive integers around the mean. Then $M_0$ fixes the basis
vectors with $j\in\mathcal I$. The transposes of the other,
orthogonal basis vectors form $C$; their entries are $0,1,-1$.
Lemma~\ref{lem:swaps-matrix} gives
\[
 c=D\Pr[W\notin\mathcal I].
\]
The standard deviation of $W$ is $\sqrt{Q\cdot\frac{r-1}{2r}\cdot\frac{r+1}{2r}}=\Theta(\sqrt Q)$ for $r\ge2$.
Taking $R$ much larger than $\sqrt Q$ therefore puts most basis vectors inside $\mathcal I$, so only a small fraction of them need to be measured.
We can still take $R$ much smaller than $Q$, so that Taylor recovery needs only $J=1+\lfloor3nR/Q\rfloor=o(n)$ coefficients per point.

\paragraph{Returning to the original variables.}
Choose variable schedules $a_v\in[r]^Q$ and factor schedules
$d_e\in[r]^Q$, so that occurrence $(e,v)$ receives label
$u(e,v)=((a_v(t),d_e(t)))_{t\in[Q]}$. With $M$ as in
\eqref{eq:overview-incidences}, Lemma~\ref{lem:incidence-labels}
gives groups of size at most $2n/r$ and at
most $h\le2M/D$ occurrences per label, provided
\[
 12D^2\le M,\qquad 2Qr^2\le n.
\]
Place occurrences sharing a label into different copies of
$\mathcal U$. In each used coordinate put its original variable
$X_v$, and put zero in unused coordinates. These copies introduce
no new Boolean variables. Apply \eqref{eq:outline-auxiliary-row}
in the copy containing each occurrence. At zero it restores $X_v$;
at $t$, its unmeasured terms use only variables in the factor's
assigned group $\{v':a_{v'}(t)=d_e(t)\}$.

Split each of the $hc$ auxiliary measurements into its contributions
from the $r$ actual variable groups. Fixing these separate values
allows independent summation within groups; fixing only their total
would still couple the groups. Since $hD\le2M\le6n$, the resulting
matrices have at most
\[
 rhc\le rh\frac{D\tau}{3r}\le2\tau n
\]
rows in total at each point. Each row has integral coefficients
with absolute sum at most $D$: substitution and restriction cannot
increase the absolute sum of the $D$ entries of an auxiliary row.
This supplies the partitions, factor assignments, and measurement matrices with parameters $(r,\tau,D)$ from
Section~\ref{subsec:preview-separation}. The full transfer and its
kernel condition are proved in Section~\ref{subsec:embedding}.

\paragraph{Parameters and construction cost.}
Put $\lambda=\lceil\log_2 n\rceil$ and
$\nu=\lceil\log_2\lambda\rceil$, and choose
\begin{equation}\label{eq:overview-explicit-parameters}
 \begin{aligned}
 r&=\left\lfloor\frac{\sqrt\lambda}{\nu^2}\right\rfloor,
 &Q&=\left\lfloor\frac{\lambda}{128\nu}\right\rfloor,
 &R&=4\lceil\sqrt\lambda\rceil,\\
 \tau&=\frac1{r\lambda},&D&=r^{2Q}.
 \end{aligned}
\end{equation}
These are precisely the parameters of
Lemma~\ref{lem:incidence-parameters} with $\beta=1$.
Apply that lemma with $b=w=3$: each variable occurs in at most three
clauses and each clause has at most three literals, so
$b,w\le\lambda$, and its hypotheses hold with $n\le M\le3n$ for
sufficiently large $n$. The lemma gives an interval
\[
 \mathcal I=\{A_0,\ldots,A_0+R-1\},\qquad
 A_0=\left\lfloor\frac{Q(r-1)}{2r}-\frac{R-1}{2}\right\rfloor,
\]
and states that these parameters satisfy the numerical hypotheses of
Lemma~\ref{lem:layout}. With $b=w=3$, these hypotheses are the two
balancing inequalities $2(b+w)D^2=12D^2\le M$ and $2Qr^2\le n$, and
$\Pr[W\notin\mathcal I]\le\tau/(rb)=\tau/(3r)$.
It also gives $D\le2^{\lambda/128}\le(2n)^{1/128}$ and the
orders in \eqref{eq:overview-parameter-bounds}.
Lemma~\ref{lem:layout}, including the coefficient estimates of
Lemma~\ref{lem:row-coefficients}, therefore constructs these forms, partitions, factor assignments,
and measurement matrices in polynomial time. The parameter checks are exact;
only finitely many sizes fail them and are handled by enumeration.
This completes the construction needed for
\eqref{eq:overview-bounded-time}.


\subsection{A weaker algorithm for general \texorpdfstring{$\#k$-SAT}{\#k-SAT}}
\label{subsec:preview-general-sat}
The algorithm above already gives a subexponential algorithm for
general $\#k$-SAT: after disjoint sparsification, we reduce each
instance directly to a $3$-CNF in which every variable occurs at
most three times. Both steps preserve the count, as follows.

Fix $k\ge2$. We use the counting-preserving version~\cite[Appendix~A]{DHMTW}
of sparsification~\cite{IPZ}, with quantitative bounds
from~\cite[Section~6]{CIP} and
\cite[Lemma~5.14 and equation~(5.4)]{CalabroThesis}.
As recorded in Lemma~\ref{lem:sparsify}, this expresses the input count as the sum of
the counts of at most
$\operatorname{poly}_k(n,\sigma^{-1})2^{\sigma n}$ formulas $G$.
Each $G$ has the same $n$ declared variables, with at most
\[
 B=O_k\!\left((\sigma^{-1}\log(2/\sigma))^{k-1}\right)
\]
occurrences per variable, and hence at most $Bn$ nonempty clauses.
An empty clause gives count zero immediately.

First turn $G$ into a $3$-CNF without changing its count. Keep
clauses of size at most three. For each longer clause, compute its
disjunction with binary OR gates and require the final output to
be one. For a gate $z=u\vee v$, impose
\[
 (\neg u\vee z)\ \wedge\ (\neg v\vee z)\ \wedge\
 (u\vee v\vee\neg z).
\]
Here $u,v$ are input literals or earlier gate outputs, and $z$ is
fresh. These clauses enforce equality with the OR, so the gate
values are uniquely determined by the original assignment.
Consequently, a satisfying assignment of $G$ has exactly one
extension to the new formula, and an assignment not satisfying
$G$ has none. This uses $O(k)$ gates per clause.

Next reduce the occurrences. If a variable $x$ occurs $s\ge2$
times in this $3$-CNF, replace its occurrences by distinct copies
$x_1,\ldots,x_s$, preserving their signs, and add Tovey's implication
cycle~\cite[p.~86]{Tov84}:
\[
 (\neg x_1\vee x_2)\ \wedge\cdots\wedge\
 (\neg x_{s-1}\vee x_s)\ \wedge\ (\neg x_s\vee x_1).
\]
The cycle forces all copies to agree. Each copy appears in one
clause inherited from the $3$-CNF and two cycle clauses, hence
in exactly three clauses. Variables with at most one occurrence
are kept unchanged, including unused declared variables.
This standard variable-copy reduction~\cite[Proposition~4.7]{Gol08}
therefore also gives a bijection between satisfying assignments.
The final formula $\widehat G$ has $\#\widehat G=\#G$ and
\[
 n\le N:=\#\{\text{variables of }\widehat G\}=O_k(Bn).
\]
All its clauses have size at most three, including the two-literal
cycle clauses, and every variable occurs at most three times.
Thus it is exactly an instance of the setting of this overview.

Apply \eqref{eq:overview-bounded-time} to each $\widehat G$ and add
the counts. When $\sigma^{-1}$ is bounded by a fixed power of
$\log n$, so is $B$, and the resulting time is
\[
 (n+m)^{O_k(1)}
 2^{O_k\!\left(\sigma n+
          Bn(\log\log n)^2/\sqrt{\log n}\right)}.
\]
For example, choosing $\sigma$ within a constant factor of
$(\log n)^{-1/(2k)}$ yields
\[
 (n+m)^{O_k(1)}
 2^{O_k\!\left(n(\log\log n)^{k+1}/(\log n)^{1/(2k)}\right)}
   = (n+m)^{O_k(1)}2^{o(n)}
\]
with polynomial space, processing the sparsification tree and
the reduced instances one at a time. The parameter $\sigma$ may
be chosen rational.

\subsection{A faster algorithm for \texorpdfstring{$\#k$-SAT}{\#k-SAT} via counting solutions to sparse low-degree polynomial equations over \texorpdfstring{$\mathbb F_2$}{F\_2}}
\label{subsec:preview-full-case}

The running-time exponent in Section~\ref{subsec:preview-general-sat}
is $O_k(\sigma n+B\rho(n))$, where
$\rho(n)=n(\log\log n)^2/\sqrt{\log n}$. If sparsification gave
$B=O_k(1)$ at $\sigma=\Theta(\rho(n)/n)$, we would already have a
$2^{O_k(\rho(n))}$-time algorithm for general $\#k$-SAT up to
polynomial factors: copying would use only $O_k(n)$ variables per
leaf, and the number of leaves would fit within the same bound.
At this choice of $\sigma$, however, the guarantee is only
$B=(\log n)^{O_k(1)}$.

We obtain the $2^{O_k(\rho(n))}$ bound in two steps. First, the counting
algorithm of this overview extends to systems of $O(n)$ polynomial
equations over $\mathbb F_2$, each of degree at most $k$ with
polylogarithmically many monomials
(Theorem~\ref{thm:overview-sparse-polynomials}). Second, even if a
leaf $G$ has $Bn$ clauses, we can construct such a system on the
original $n$ variables whose common zeros are exactly the satisfying
assignments of $G$. We illustrate the reduction first with a
self-contained randomized construction, then explain the counting
algorithm.

\medskip\noindent
\textbf{A randomized reduction to sparse polynomial equations.}
Write $v_e(x)=1$ if $x$ violates clause $e$, and $0$ otherwise. A
\emph{parity test} on a set $S$ of clauses is the polynomial equation
$\sum_{e\in S}v_e(x)=0$ over $\mathbb F_2$: it accepts if $x$ violates
an even number of clauses in $S$. Every satisfying assignment passes,
so we need only catch every nonsatisfying one. If $G$ has $m_G\le n$
clauses, one test per clause suffices. Suppose $m_G>n$.

A short test can easily miss a single violated clause, but all
assignments that violate only that clause give identical answers.
Thus there are only $m_G$ singleton patterns to catch. Use $8n$
independent tests. For each, draw $L=2B\lceil\log_2(2m_G)\rceil$
clauses uniformly with replacement. For each draw, independently add
its violation bit to the parity sum with probability $1/2$; two
included copies cancel. Since $m_G\le Bn$, each clause is drawn
at least logarithmically often in expectation overall. For a fixed
set of $w\ge1$ violated clauses, let $u=Lw/m_G$ be the expected number
of hits per test. If the sample hits the set, the coin on the first
hit makes the parity odd with probability $1/2$, even with all other
coins fixed. Thus the rejection probability is
\[
 p_w=\tfrac12\bigl(1-(1-w/m_G)^L\bigr)\ge\tfrac12(1-e^{-u}).
\]

We count \emph{sets} for few violations and \emph{assignments} for many:
\begin{itemize}
\item If $u\le1$, then $p_w\ge u/4$. There are at most $m_G^w$ sets
of size $w$, so the chance that at least one escapes all tests is
at most $m_G^w e^{-8np_w}\le m_G^w e^{-2nLw/m_G}\le16^{-w}$.
Summing over $w\ge1$ gives failure probability at most $1/15$.
\item If $u>1$, then $p_w>1/4$. A union bound over the $2^n$
assignments bounds the chance that any assignment in this case
escapes by $2^n(3/4)^{8n}<1/4$.
\end{itemize}
Thus with probability greater than $2/3$, the common zeros are
exactly the satisfying assignments of $G$. Each $v_e$ is a product
of at most $k$ factors of the form $x_j$ or $1+x_j$, so each test has degree at
most $k$ and at most $2^kL$ monomials, polylogarithmic in $n$.
Section~\ref{sec:hashes} constructs short tests with the same
solution set deterministically.

\medskip\noindent
\textbf{Counting the common zeros.}
Recall from Section~\ref{subsec:preview-encoding} that $\Phi_e(x)$ is
$1$ if $x$ satisfies clause $e$ and $0$ otherwise. What should replace
this factor when we test a whole set $S_i$ of clauses at once? First
count how many clauses in \emph{this test} are violated, using ordinary
integer addition:
\[
 \Lambda_i(x)=\sum_{e\in S_i}v_e(x)=\sum_{e\in S_i}\bigl(1-\Phi_e(x)\bigr).
\]
The test passes exactly when $\Lambda_i(x)$ is even. Choose $\ell=n+1$, so
that a count of at most $2^n$ is determined by its residue modulo
$2^\ell$. The complemented Beigel--Tarui polynomial $p_\ell$
(Lemma~\ref{lem:rational-idempotent}) is $1$ modulo $2^\ell$ on even
integers and $0$ on odd integers. For example, violating two clauses
in $S_i$ still passes the test: $p_\ell(2)\equiv1\pmod{2^\ell}$.
Writing $m_{\mathrm{test}}$ for the number of tests, our counting expression becomes
\[
 \underbrace{\sum_x\prod_{e\in G}\Phi_e(x)}_{\text{one factor per clause}}
 \qquad\longrightarrow\qquad
 \underbrace{\sum_x\prod_{i=1}^{m_{\mathrm{test}}} p_\ell(\Lambda_i(x))}_{\text{one factor per test}}.
\]
The left side is $\#G$. The least nonnegative residue of the right
side modulo $2^\ell$ is the number of assignments passing every test.
When the tests catch every nonsatisfying assignment, this is exactly
$\#G$. The same encoding works for
any sparse polynomial equation $P_i(x)=0$ over $\mathbb F_2$: take
$L_i(x)$ to be the integer sum of its listed monomials. Again the
equation holds exactly when this sum is even.

To reuse the Taylor step, we need a product of input degree $O(n)$.
Multiplying the $m_{\mathrm{test}}$ test factors directly gives only the bound
$O(km_{\mathrm{test}}\ell)$, potentially quadratic in $n$. The key is to truncate
the product \emph{across all tests together}. Section~\ref{subsec:rational-lift}
writes $p_\ell(y)$ as a sum of terms $e_j(y)h(y)^j$, where
$e_j$ has degree at most one and $h(y)=y(y-1)$ is even at every
integer. A product term containing $\ell$ or more copies of $h$ is
therefore zero modulo $2^\ell$ on Boolean assignments, even if those
copies come from different tests. We keep just the terms with fewer
than $\ell$ copies in total, tracking that number as we multiply.
Each retained term has at most $m_{\mathrm{test}}$ linear factors and $\ell-1$
quadratic factors in the $L_i$'s. Since $\deg L_i\le k$, its input
degree is at most $k(m_{\mathrm{test}}+2(\ell-1))=O(n)$.

For the replacement forms, what matters is which variables each factor
reads. Where Section~\ref{subsec:preview-deformation} used one form
for each variable in a clause, we now use one for each variable in a
test, reusing it throughout that test. Each sparse equation reads only
polylogarithmically many variables. Branching on the values of the
$O(n/\sqrt{\log n})$ most frequent variables leaves each remaining
variable in only polylogarithmically many tests, as required by
Section~\ref{sec:construction}. On each branch, fixing the shared
measurements again lets us sum one group at a time. Because the
truncated product has input degree $O(n)$, the degree-$R$ replacements
give $\deg_T H=O(nR)$: the same $O(1+nR/Q)$ Taylor coefficients per
point recover the contributions at $T=0$ with $0,1,\ldots,\ell-1$
copies of $h$. Their sum modulo $2^\ell$, taken as the least nonnegative
residue, is the exact branch count. Adding the branch counts gives the
number of common zeros.

\begin{theorem}[Counting common zeros of sparse polynomials]
\label{thm:overview-sparse-polynomials}
Fix integers $k\ge2$ and $c\ge1$. Let $P_1,\ldots,P_{m_{\mathrm{test}}}$ be polynomials
over $\mathbb F_2$ in $n\ge20$ Boolean variables, with $m_{\mathrm{test}}\le cn$.
Suppose each $P_i$ is given as a list of at most $(\log n)^c$
monomials of degree at most $k$. The exact number of common zeros in
$\{0,1\}^n$ can be computed deterministically in
$2^{O_{k,c}(n(\log\log n)^2/\sqrt{\log n})}$ bit operations and
$n^{O_{k,c}(1)}$ space.
\end{theorem}
Section~\ref{sec:setup} proves the theorem.

For the final system constructed for $G$, the number of common zeros is $\#G$.
Section~\ref{sec:hashes} uses only $2^{(\log n)^{O_k(1)}}$ calls
to the counter per leaf, which preserves the running-time exponent in
Theorem~\ref{thm:overview-sparse-polynomials}. Finally, choosing
$\sigma$ of order $(\log\log n)^2/\sqrt{\log n}$ in
Section~\ref{subsec:general-reduction} keeps $B$ polylogarithmic
and the number of leaves within $2^{O_k(n(\log\log n)^2/\sqrt{\log n})}$.
Their satisfying assignments are disjoint, so adding their counts
gives the bound in Proposition~\ref{prop:interpolation-bound}.

